\documentclass[11pt,a4paper]{article}

\usepackage{geometry}
\geometry{margin=1in}
\usepackage{amsmath,amssymb,amsfonts}
\usepackage{booktabs}
\usepackage{array}
\usepackage{longtable}
\usepackage{multirow}
\usepackage{siunitx}
\usepackage{graphicx}
\usepackage{xcolor}
\usepackage{hyperref}
\usepackage{enumitem}
\usepackage{listings}
\usepackage{float}
\usepackage{caption}
\usepackage{fontspec}
\usepackage{xeCJK}

\IfFontExistsTF{Noto Serif CJK TC}{\setCJKmainfont{Noto Serif CJK TC}}{%
\IfFontExistsTF{Source Han Serif TC}{\setCJKmainfont{Source Han Serif TC}}{%
\IfFontExistsTF{PingFang TC}{\setCJKmainfont{PingFang TC}}{}}}

\hypersetup{
    colorlinks=true,
    linkcolor=blue!60!black,
    urlcolor=blue!60!black,
    citecolor=blue!60!black
}

\lstset{
    basicstyle=\ttfamily\small,
    breaklines=true,
    frame=single,
    columns=fullflexible,
    backgroundcolor=\color{gray!5}
}

\title{SpMDV Sparse Matrix--Dense Vector Multiplication\\Architecture Exploration Report}
\author{P. Cheng}
\date{\today}

\begin{document}
\maketitle

\begin{abstract}
This report summarizes the design exploration for a sparse matrix--dense vector multiplication (SpMDV) accelerator. The target computation is a bank-balanced sparse matrix multiplied by dense vector batches. The main architectural knobs are matrix-bank parallelism and dense-vector feature parallelism. We derive a cycle and area model, describe the implemented 64-MAC architecture, explain the practical optimizations, and summarize current synthesis and simulation results. The current 64-MAC design reduces area-time product (AT) compared with the 1-MAC reference, but the improvement is limited by SRAM banking overhead and output pipeline overhead. The next recommended experiments are a 32-MAC design and optimized 64-MAC variants.
\end{abstract}

\section{Problem Statement}

The target module implements sparse matrix--dense vector multiplication:
\begin{equation}
    y = Wx + b.
\end{equation}

The matrix $W$ is a $256 \times 256$ sparse matrix stored using Bank-Balanced Sparsity (BBS). Each matrix row is partitioned into four column banks. Each bank contains 12 non-zero values per row. Therefore, each row contains
\begin{equation}
    K = 4 \times 12 = 48
\end{equation}
non-zero values.

The testbench repeatedly provides batches of dense vectors. In the current task setting, each batch contains 16 dense vectors, and the full testbench contains two such batches, for a total of 32 dense vectors. The dense vector dimension is 256.

\subsection{Data Formats}

The fixed-point formats must remain unchanged across all architectures:

\begin{table}[H]
\centering
\caption{Fixed-point format requirement.}
\begin{tabular}{lll}
\toprule
Data & Width & Format \\
\midrule
Weight & 8-bit signed & $S2.5$ \\
Dense vector & 8-bit signed & $S1.6$ \\
Product & 16-bit signed & $S4.11$ \\
Bias & 8-bit signed & $S0.7$ \\
Output & 22-bit signed & $S10.11$ \\
\bottomrule
\end{tabular}
\end{table}

The product format is preserved as
\begin{equation}
    S2.5 \times S1.6 = S4.11.
\end{equation}
Bias is converted into output format by sign extension and appending four fractional zeros:
\begin{equation}
    b_{S10.11} = \text{sign\_extend}(b_{S0.7}) \ll 4.
\end{equation}
In Verilog-style notation:
\begin{lstlisting}[language=Verilog]
bias_to_s10_11 = $signed({{10{bias_byte[7]}}, bias_byte, 4'b0000});
\end{lstlisting}

\subsection{Input Protocol}

The external input protocol is fixed and should not be changed.

\paragraph{Matrix initialization stream.}
The testbench sends initialization data through \texttt{raw\_input[7:0]} with \texttt{w\_input\_valid}:
\begin{enumerate}[nosep]
    \item all weight values: $256 \times 48 = 12288$ bytes,
    \item all position values: $256 \times 48 = 12288$ bytes,
    \item all bias values: $256$ bytes.
\end{enumerate}
The total initialization stream is
\begin{equation}
    12288 + 12288 + 256 = 24832 \text{ bytes}.
\end{equation}

For weight and position, the stream order is row, group, bank:
\begin{equation}
    (r,g,b) = (0,0,0), (0,0,1), (0,0,2), (0,0,3), (0,1,0), \ldots.
\end{equation}

\paragraph{Dense vector stream.}
Each batch contains 16 dense vectors, provided in feature-major order:
\begin{equation}
    x_0[0..255], x_1[0..255], \ldots, x_{15}[0..255].
\end{equation}
Thus each batch requires
\begin{equation}
    16 \times 256 = 4096
\end{equation}
cycles of vector loading when the input bandwidth is one byte per cycle.

\paragraph{Output order.}
The output is expected in feature-major order:
\begin{equation}
    y_0[0..255], y_1[0..255], \ldots, y_{15}[0..255].
\end{equation}
This output order is a major reason why a result buffer is required for row-stationary parallel computation.

\section{Theory Model Building}

\subsection{Fixed Problem Parameters}

We define the following fixed parameters:
\begin{align*}
    R &= 256 && \text{number of rows},\\
    C &= 256 && \text{dense vector length},\\
    B &= 4   && \text{BBS matrix banks},\\
    G &= 12  && \text{non-zero groups per bank per row},\\
    F &= 16  && \text{dense vectors per batch}.
\end{align*}

Each row has
\begin{equation}
    K = B G = 48
\end{equation}
non-zero values. The total number of scalar multiply operations per 16-vector batch is
\begin{equation}
    N_{\mathrm{MAC}} = R F B G = 256 \cdot 16 \cdot 4 \cdot 12 = 196608.
\end{equation}

\subsection{Parallelism Parameters}

The two architectural knobs are:
\begin{align*}
    p &= \text{number of matrix banks processed per cycle},\\
    n &= \text{number of dense vector features processed per cycle}.
\end{align*}

The number of physical MAC lanes is
\begin{equation}
    m = p n.
\end{equation}

The most useful design line is usually to first maximize matrix-bank parallelism, because BBS naturally provides four independent banks. Therefore, the primary exploration points are:
\begin{equation}
    (p,n) \in \{(1,1), (4,4), (4,8), (4,16)\},
\end{equation}
corresponding to 1-MAC, 16-MAC, 32-MAC, and 64-MAC.

\subsection{Compute Cycle Model}

A row contains $G$ groups. For each group, there are $B$ matrix banks and $F$ dense-vector features. If we process $p$ banks and $n$ features per issue, the ideal compute cycles per batch are
\begin{equation}
    T_{\mathrm{compute}}(p,n)
    = R G \left\lceil \frac{B}{p} \right\rceil
          \left\lceil \frac{F}{n} \right\rceil.
\end{equation}

When $p$ divides $B$ and $n$ divides $F$,
\begin{equation}
    T_{\mathrm{compute}}(p,n)
    = \frac{R F B G}{pn}
    = \frac{196608}{pn}.
\end{equation}

For the main design points:
\begin{table}[H]
\centering
\caption{Ideal compute cycle model per 16-vector batch.}
\begin{tabular}{rrrrr}
\toprule
Architecture & $p$ & $n$ & MAC lanes $m=pn$ & $T_{\mathrm{compute}}$ \\
\midrule
1-MAC & 1 & 1 & 1 & 196608 \\
16-MAC & 4 & 4 & 16 & 12288 \\
32-MAC & 4 & 8 & 32 & 6144 \\
64-MAC & 4 & 16 & 64 & 3072 \\
\bottomrule
\end{tabular}
\end{table}

\subsection{Non-compute Cycle Components}

The total measured runtime is not only compute cycles. It also includes vector loading, result buffering, output drain, and pipeline overhead.

Dense-vector loading costs
\begin{equation}
    T_{\mathrm{vec-load}} = F C = 4096
\end{equation}
cycles per batch, assuming one byte per cycle input.

If the output is not pipelined and each result requires request, wait, and send states, then
\begin{equation}
    T_{\mathrm{output,old}} = 3 R F = 12288
\end{equation}
cycles per batch. With an output read pipeline, after fill, the output can approach
\begin{equation}
    T_{\mathrm{output,new}} \approx R F + L_{\mathrm{out}} = 4096 + L_{\mathrm{out}},
\end{equation}
where $L_{\mathrm{out}}$ is a small pipeline fill/drain constant.

\subsection{SRAM Area Model}

The measured SRAM macro areas are:
\begin{align*}
    A_{256}  &= 20810,\\
    A_{512}  &= 30948,\\
    A_{4096} &= 131907.
\end{align*}

For a true 4-bank BBS matrix layout, matrix storage requires:
\begin{equation}
    4 \text{ weight SRAMs} + 4 \text{ position SRAMs} = 8 A_{4096}.
\end{equation}

For a $4 \times n$ architecture:
\begin{itemize}[nosep]
    \item vector buffer requires $4n$ SRAMs of size $256 \times 8$,
    \item result buffer requires $3n$ SRAMs of size $256 \times 8$ if only one $n$-feature block is buffered,
    \item bias requires one $256 \times 8$ SRAM.
\end{itemize}

Therefore,
\begin{equation}
    A_{\mathrm{SRAM}}(n) = 8A_{4096} + (1 + 7n)A_{256}.
\end{equation}

For the main 4-bank designs:
\begin{table}[H]
\centering
\caption{SRAM area model for $4 \times n$ architectures.}
\begin{tabular}{rrrrr}
\toprule
Architecture & $n$ & Vector SRAMs & Result SRAMs & $A_{\mathrm{SRAM}}$ \\
\midrule
16-MAC & 4  & 16 & 12 & 1,658,746 \\
32-MAC & 8  & 32 & 24 & 2,241,426 \\
64-MAC & 16 & 64 & 48 & 3,406,786 \\
\bottomrule
\end{tabular}
\end{table}

\subsection{Area-Time Product Model}

We use
\begin{equation}
    AT = A \cdot T
\end{equation}
where $A$ is synthesized area and $T$ is measured cycle count or real runtime. If clock periods differ substantially, the more realistic metric is
\begin{equation}
    AT_{\mathrm{real}} = A \cdot T_{\mathrm{cycles}} \cdot T_{\mathrm{clk}}.
\end{equation}

The key trade-off is:
\begin{itemize}[nosep]
    \item increasing $pn$ reduces compute cycles,
    \item increasing $n$ increases vector SRAM and result SRAM area,
    \item for high parallelism, fixed overheads such as vector loading and output drain limit the benefit of reducing compute cycles.
\end{itemize}

\section{Practical Optimizations Adopted}

\subsection{True BBS Matrix Banking}

Instead of storing all $12288$ weights across three capacity-based SRAMs, the optimized architecture stores matrix values according to their BBS bank:
\begin{align*}
    W_b[r,g] &: b \in \{0,1,2,3\},\\
    P_b[r,g] &: b \in \{0,1,2,3\}.
\end{align*}
Each bank contains $256 \times 12 = 3072$ entries and fits into one $4096 \times 8$ SRAM. The address is
\begin{equation}
    \mathrm{addr}(r,g) = 12r + g = 8r + 4r + g.
\end{equation}

\subsection{64-MAC Vector Banking}

For the 64-MAC design, the vector buffer is organized as:
\begin{equation}
    X[f,b,\ell], \quad f \in [0,15],\ b \in [0,3],\ \ell \in [0,63].
\end{equation}
Each feature and column bank has its own $256 \times 8$ SRAM. Therefore, the 64-MAC design uses
\begin{equation}
    16 \times 4 = 64
\end{equation}
vector SRAMs. This gives enough read bandwidth for 64 vector values per compute issue.

\subsection{Result Buffering}

The output is 22-bit signed $S10.11$. It is stored in three byte SRAMs per feature:
\begin{align*}
    \mathrm{low}  &= y[7:0],\\
    \mathrm{mid}  &= y[15:8],\\
    \mathrm{high} &= \{2'b00, y[21:16]\}.
\end{align*}
For 16 features, the 64-MAC design uses
\begin{equation}
    16 \times 3 = 48
\end{equation}
result SRAMs.

\subsection{Compute Pipeline}

The compute pipeline hides synchronous SRAM latency. A single group bundle has multiple stages:
\begin{enumerate}[nosep]
    \item issue four weight/position SRAM reads,
    \item wait for matrix SRAM output,
    \item issue 64 vector SRAM reads,
    \item multiply 64 pairs and reduce four products per feature,
    \item accumulate 16 partial sums into 16 accumulators.
\end{enumerate}

The key idea is that the pipeline latency is several cycles, but the issue rate is one group bundle per cycle after fill. Therefore, a row with 12 groups requires roughly
\begin{equation}
    12 + L_{\mathrm{pipe}}
\end{equation}
cycles, not $12 \times L_{\mathrm{pipe}}$.

\subsection{Optimizations for the 64-MAC Version}

The first 64-MAC implementation used the following expensive sequence:
\begin{equation*}
    \text{bias} \rightarrow \text{result SRAM} \rightarrow \text{accumulators} \rightarrow \text{compute} \rightarrow \text{result SRAM} \rightarrow \text{output}.
\end{equation*}
Three optimizations were introduced.

\paragraph{Output pipeline.}
The old output FSM used request, wait, and send states per result. This costs about $3RF=12288$ cycles per batch. The optimized version issues the next result SRAM read every cycle and outputs the result after pipeline latency. This should approach $RF=4096$ cycles per batch.

\paragraph{Accumulator bias initialization.}
Instead of pre-filling result SRAM with bias and reading it back into accumulators, the optimized version directly initializes accumulators:
\begin{equation}
    \mathrm{acc}[f] \leftarrow b[r]_{S10.11}, \quad f=0,1,\ldots,15.
\end{equation}
This removes the result-SRAM bias-fill phase and the result-SRAM read-to-accumulator phase.

\paragraph{Bias prefetch.}
While computing row $r$, the bias for row $r+1$ can be prefetched from the independent bias SRAM. This hides the bias read bubble for the next row.

\section{Experiment Results}

\subsection{Current 64-MAC Baseline Result}

The current synthesized 64-MAC design has the following area distribution:

\begin{table}[H]
\centering
\caption{Synthesized area of the current 64-MAC design.}
\begin{tabular}{lr}
\toprule
Component & Area \\
\midrule
Total area & 3,729,240.1982 \\
Combinational area reported at top level & 45,746.6268 \\
Non-combinational area & 73,721.4777 \\
Black-box SRAM area & 3,406,774.9355 \\
\bottomrule
\end{tabular}
\end{table}

The black-box SRAM area dominates:
\begin{equation}
    \frac{3,406,774.9355}{3,729,240.1982} \approx 91.35\%.
\end{equation}
This confirms that the 64-MAC architecture is memory-area dominated rather than multiplier-area dominated.

The measured runtime for two 16-vector batches is:
\begin{equation}
    T_{64,\mathrm{old}} = 45570 \text{ cycles}.
\end{equation}
Thus, per 16-vector batch:
\begin{equation}
    T_{64,\mathrm{old,batch}} \approx 22785 \text{ cycles}.
\end{equation}

The corresponding cycle-based area-time product is
\begin{equation}
    AT_{64,\mathrm{old}} = 3,729,240.1982 \times 45,570
    \approx 1.6994 \times 10^{11}.
\end{equation}

This version was observed to reduce AT by roughly 40\% compared with the original 1-MAC reference. However, the improvement is limited by SRAM area and non-compute cycle overhead.

\subsection{Cycle Breakdown of the Old 64-MAC Version}

The measured per-batch cycle count is almost exactly explained by:

\begin{table}[H]
\centering
\caption{Approximate cycle breakdown of the old 64-MAC implementation.}
\begin{tabular}{lr}
\toprule
Phase & Cycles per batch \\
\midrule
Read 16 dense vectors & 4096 \\
Bias-fill result SRAM & $256 \times 3 = 768$ \\
Row compute and result writeback & $256 \times 22 = 5632$ \\
Output with REQ/WAIT/SEND & $4096 \times 3 = 12288$ \\
\midrule
Total & 22784 \\
Measured & 22785 \\
\bottomrule
\end{tabular}
\end{table}

This shows that output drain is the largest contributor. The ideal 64-MAC compute body is only $3072$ cycles per batch, but the old output path alone consumes $12288$ cycles per batch.

\subsection{Expected Impact of the New 64-MAC Optimizations}

The optimized 64-MAC version targets the following reductions:

\begin{table}[H]
\centering
\caption{Expected cycle reductions in optimized 64-MAC implementation.}
\begin{tabular}{lrr}
\toprule
Optimization & Old cost & Expected new cost \\
\midrule
Output pipeline & $12288$ & $4096 + L_{\mathrm{out}}$ \\
Remove bias-fill result SRAM & $768$ & $0$ \\
Remove result SRAM read-to-acc & included in row overhead & mostly removed \\
Bias prefetch & row bias bubbles & hidden after first row \\
\bottomrule
\end{tabular}
\end{table}

A rough target for two 16-vector batches is:
\begin{equation}
    T_{64,\mathrm{opt}} \approx 26000 \sim 30000 \text{ cycles}.
\end{equation}
This number should be confirmed by RTL and gate-level simulation.

\section{Discussion: Why 32-MAC Still Matters}

Although the 64-MAC version improves AT, it uses a large amount of SRAM. The 64-MAC design requires 64 vector SRAMs and 48 result SRAMs. Reducing feature parallelism from 16 to 8 gives a 32-MAC design with:
\begin{align*}
    \text{vector SRAMs} &= 4 \times 8 = 32,\\
    \text{result SRAMs} &= 3 \times 8 = 24.
\end{align*}

The SRAM area decreases from
\begin{equation}
    A_{\mathrm{SRAM}}(16) = 3,406,786
\end{equation}
to
\begin{equation}
    A_{\mathrm{SRAM}}(8) = 2,241,426.
\end{equation}
The compute cycles double from 3072 to 6144 per batch, but the total runtime does not double because vector loading and output remain significant. Therefore, the 32-MAC design is a strong candidate for the best AT point.

\section{Conclusion}

The SpMDV problem is governed by two independent parallelism axes:
\begin{equation}
    \mathrm{MACs} = \mathrm{matrix\ bank\ parallelism} \times \mathrm{feature\ parallelism}.
\end{equation}

The matrix side naturally supports four-way parallelism due to BBS. Feature parallelism provides additional throughput but requires significant vector and result SRAM banking. The 64-MAC design demonstrates that aggressive parallelism can reduce AT, but the benefit is limited by SRAM area and output overhead. The synthesis results show that more than 91\% of the 64-MAC area is black-box SRAM, so future improvements should focus on memory architecture and cycle overhead rather than multiplier area.

The immediate next steps are:
\begin{enumerate}[nosep]
    \item validate the optimized 64-MAC version with output pipeline, accumulator bias initialization, and bias prefetch;
    \item synthesize and simulate the 32-MAC $4 \times 8$ architecture;
    \item compare 1-MAC, 16-MAC, 32-MAC, old 64-MAC, and optimized 64-MAC using area, cycle count, clock period, and AT;
    \item use gate-level timing results to evaluate whether higher parallelism increases clock period enough to offset cycle reductions.
\end{enumerate}

The current best interpretation is that 64-MAC is a useful extreme-throughput reference, while 32-MAC is likely the most important design point for area-time efficiency.

\appendix
\section{Recommended Result Table Template}

After all experiments are complete, fill the following table:

\begin{table}[H]
\centering
\caption{Final experiment comparison template.}
\begin{tabular}{lrrrrrr}
\toprule
Version & $p$ & $n$ & MACs & Area & Cycles & AT \\
\midrule
1-MAC reference & 1 & 1 & 1 & -- & -- & -- \\
16-MAC $4\times4$ & 4 & 4 & 16 & -- & -- & -- \\
32-MAC $4\times8$ & 4 & 8 & 32 & -- & -- & -- \\
64-MAC old $4\times16$ & 4 & 16 & 64 & 3,729,240 & 45,570 & $1.6994\times10^{11}$ \\
64-MAC optimized $4\times16$ & 4 & 16 & 64 & -- & -- & -- \\
\bottomrule
\end{tabular}
\end{table}

\section{Important Implementation Invariants}

All optimized versions should preserve:
\begin{enumerate}[nosep]
    \item the original top-level I/O interface;
    \item the original request/valid protocol;
    \item the external stream order: weights, positions, biases, then dense vectors;
    \item all fixed-point formats and sign-extension rules;
    \item synchronous SRAM latency handling;
    \item feature-major output order.
\end{enumerate}

\end{document}
